Surrogate-assisted optimization of small tailless aircraft is usually judged against its initial sample, which does not show what the surrogate contributes. This paper benchmarks eight search strategies under an equal budget on the conceptual design of a 1.5\,m, 8\,kg electric blended-wing-body UAV with 14 variables and a vortex-lattice and strip-drag evaluator. Every strategy receives 287 evaluations beyond a shared pool of 9,800 designs, with 10 to 20 seeds. Trust-region Bayesian optimization (TuRBO) gives the longest calm-air range, $\NvMeanTurbo \pm \NvSdTurbo$\,km. A neural-ensemble active-learning loop weighted by a feasibility classifier reaches $\NvMeanAL \pm \NvSdAL$\,km as originally implemented, no better than an adaptive local search without a surrogate ($\NvMeanAdapt \pm \NvSdAdapt$\,km). After a repair that stops the loop from proposing designs that violate the washout ordering it reaches $\NvMeanALR \pm \NvSdALR$\,km, \NvPctTurbovALR\% below TuRBO (95\% interval $\NvCITurbovALR$\,km). Without the classifier weight the loop finds no feasible design in 20 runs, whereas the same weight does not improve TuRBO. An audit of the code against the text corrected two baseline implementations and withdrew a monotonicity claim for the planform loft. Euler CFD on \NcalN\ designs shows that the lattice underpredicts the static margin by a nearly constant \NcalA, which moves the effective upper limit from 0.25 to \NcalCapEquiv\ and places the best design of a minority of runs beyond it without changing the ranking of the strategies. The ordering persists for pools of 288 to 3{,}000 designs, and a pool of 1{,}000 designs costs TuRBO \NpTurboLossB\ to \NpIndTurboVsFull\,km. The conclusions hold for this problem and evaluator.
